% © 2026 Ing. David Jaroš — CC BY-NC-ND 4.0
%
% This work is licensed under a Creative Commons Attribution-NonCommercial-NoDerivatives
% 4.0 International License (CC BY-NC-ND 4.0).
% See LICENSE.md for full license text.

% UBT-AI-PROVENANCE-BEGIN
% schema: ubt-ai-provenance/v1
% tier: C_working
% ai_assistance: disclosed
% human_review: risk-based
% editorial_responsibility: Ing. David Jaroš
% policy: ../../AI_PROVENANCE.md
% notice: Working material; exhaustive human review is not claimed.
% UBT-AI-PROVENANCE-END

\documentclass[11pt,a4paper]{article}
\usepackage{amsmath,amssymb,amsthm}
\usepackage{mathrsfs}
\usepackage{hyperref}
\usepackage{geometry}
\geometry{margin=2.5cm}

\newtheorem{definition}{Definition}
\newtheorem{proposition}{Proposition}
\newtheorem{remark}{Remark}
\newtheorem{conjecture}{Conjecture}
\newtheorem{gap}{Open Gap}

% Proof-status commands
\newcommand{\proved}[1]{\textbf{[PROVED]} #1}
\newcommand{\heuristic}[1]{\textbf{[HEURISTIC]} #1}
\newcommand{\speculative}[1]{\textbf{[SPECULATIVE]} #1}

\title{Candidate Self-Adjoint Operator for the UBT Hilbert--P\'olya Programme:\\
Operator Definition and Hilbert Space Framework}
\author{Ing.\ David Jaro\v{s}\\
\small Unified Biquaternion Theory Research\\
\small \texttt{research\_tracks/rh\_operator/}}
\date{May 2026}

\begin{document}
\maketitle

\begin{abstract}
We construct an explicit candidate self-adjoint operator $\hat{H}_\psi$ arising
from the $\psi$-sector of the Unified Biquaternion Theory (UBT), to be compared
against the spectrum of the Riemann zeta function within the Hilbert--P\'olya
programme.  We define the Hilbert space, domain, and boundary conditions,
state a candidate effective potential, and summarise the open gaps that must be
closed before any connection to the Riemann Hypothesis (RH) can be claimed.

\medskip
\noindent\textbf{Proof-status key:} \proved{} established rigorously;
\heuristic{} plausible but not proved; \speculative{} conjectural.
\end{abstract}

\tableofcontents

\bigskip
\noindent\textbf{Critical warning.}  This document \emph{does not} claim or
imply a proof of the Riemann Hypothesis.  Any construction that takes the
non-trivial zeros $\tfrac12+i\gamma_n$ as input is circular.  All claims are
explicitly labelled by proof status.

%------------------------------------------------------------
\section{Background: The Hilbert--P\'olya Programme}
%------------------------------------------------------------

\subsection{Conjecture}

The Hilbert--P\'olya conjecture asserts the existence of a self-adjoint operator
$T$ on a Hilbert space $\mathscr{H}$ such that
\[
  \operatorname{spec}(T) = \bigl\{\,\gamma_n \in \mathbb{R} :
  \zeta(\tfrac12 + i\gamma_n) = 0,\; \gamma_n > 0\,\bigr\}.
\]
Under RH the zeros lie on the critical line, so the spectrum is real, which is
automatic for a self-adjoint operator.

\subsection{Known candidate operators}

\begin{enumerate}
\item \textbf{Berry--Keating} \cite{BK99}: semiclassical $H = xp$ on the
      half-line.  The classical periodic orbits produce a prime-counting
      contribution, but the quantisation is not fully defined.
\item \textbf{Connes} \cite{C99}: adelic trace formula; cokernel of a
      multiplication operator on $L^2(\mathbb{A}_\mathbb{Q}/\mathbb{Q}^*)$.
\item \textbf{Bost--Connes} \cite{BC95}: $C^*$-algebraic system whose
      partition function is $\zeta(s)$ at inverse temperature $\beta = s$.
\item \textbf{Schr\"odinger--zeta operators} (various): $-d^2/dx^2 + V(x)$
      with $V$ chosen to mimic zero spacings.
\end{enumerate}

%------------------------------------------------------------
\section{UBT \texorpdfstring{$\psi$}{psi}-Sector Framework}
%------------------------------------------------------------

\subsection{Fundamental field}

The UBT fundamental field is
\[
  \Theta(q,\tau), \qquad q \in \mathbb{C}\otimes\mathbb{H}, \quad
  \tau = t + i\psi \in \mathbb{C}.
\]
The imaginary time component $\psi$ parameterises a compact direction of
circumference $L_\psi$ (to be determined from the UBT moduli).  Integrating the
UBT Lagrangian over all spatial biquaternion directions $q$ yields an effective
1-dimensional theory in the $\psi$ variable.

\subsection{Hilbert space}
\label{sec:hilbert}

\begin{definition}[Hilbert space]
  \label{def:hilbert}
  \[
    \mathscr{H}_\psi := L^2\bigl(S^1_{L_\psi},\, d\psi\bigr)
  \]
  where $S^1_{L_\psi} = \mathbb{R}/(L_\psi\mathbb{Z})$ is the circle of
  circumference $L_\psi$.  The inner product is
  \[
    \langle f, g \rangle
    = \int_0^{L_\psi} \overline{f(\psi)}\, g(\psi)\, d\psi.
  \]
\end{definition}

\textbf{Proof status}: \proved{The space $L^2(S^1_{L_\psi})$ is a standard
separable Hilbert space by the theory of square-integrable functions on a
compact Riemannian manifold.}  The identification $L_\psi = 2\pi$ is
\heuristic{motivated by dimensional analysis; an independent derivation from
UBT moduli is an open gap (Gap~\ref{gap:Lpsi}).}

%------------------------------------------------------------
\section{Candidate Operator}
%------------------------------------------------------------

\subsection{Formal definition}

\begin{definition}[Candidate UBT $\psi$-operator]
  \label{def:operator}
  Let $D_0 = C^\infty_\mathrm{per}(S^1_{L_\psi})$ be the space of smooth
  $L_\psi$-periodic functions.  Define
  \[
    \hat{H}_\psi := -\frac{d^2}{d\psi^2} + V_\mathrm{eff}(\psi)
  \]
  on the domain $D_0 \subset \mathscr{H}_\psi$, with periodic boundary
  conditions
  \[
    f(0) = f(L_\psi), \qquad f'(0) = f'(L_\psi).
  \]
\end{definition}

\subsection{Candidate effective potential}
\label{sec:Veff}

From the UBT curvature of the $\psi$-direction, the effective potential is
expected to take the form
\[
  V_\mathrm{eff}(\psi) = V_0 + \kappa\,\langle\Theta^\dagger\Theta\rangle_q
  \big|_\psi,
\]
where the angle brackets denote integration over the spatial biquaternion
directions and $\kappa$ is the UBT coupling constant.  In the prime-stability
context (cf.\ canonical/alpha/) the discrete analogue is
\[
  V_\mathrm{eff}(n) = n^2 - B\, n\ln n, \qquad
  B \approx 46 \text{ (empirically)}.
\]

\textbf{Proof status}: The form $n^2 - Bn\ln n$ is
\heuristic{motivated by a one-loop renormalisation group argument (see
\texttt{rg\_nlogn/} track); a full derivation of $V_\mathrm{eff}(\psi)$
from the UBT Lagrangian is an open gap (Gap~\ref{gap:Veff}).}

\subsection{Free ($V_\mathrm{eff}=0$) case}

\begin{proposition}[Free spectrum]
  When $V_\mathrm{eff} = 0$, the operator $\hat{H}_\psi = -d^2/d\psi^2$
  with periodic boundary conditions on $S^1_{L_\psi}$ is essentially
  self-adjoint on $D_0$ and has spectrum
  \[
    \operatorname{spec}(\hat{H}_\psi^{(0)})
    = \Bigl\{\, \lambda_n = \Bigl(\frac{2\pi n}{L_\psi}\Bigr)^2 :
    n \in \mathbb{Z} \,\Bigr\}.
  \]
\end{proposition}

\textbf{Proof}: \proved{Standard Fourier analysis on $L^2(S^1_{L_\psi})$.
The eigenfunctions are $e_n(\psi) = L_\psi^{-1/2} e^{2\pi i n\psi/L_\psi}$,
which form a complete orthonormal basis.}

%------------------------------------------------------------
\section{Sturm--Liouville Form}
%------------------------------------------------------------

The operator $\hat{H}_\psi = -d^2/d\psi^2 + V_\mathrm{eff}(\psi)$ is in
standard Sturm--Liouville form
\[
  \mathcal{L}f = -\frac{d}{d\psi}\!\left[p(\psi)\frac{df}{d\psi}\right]
  + q(\psi)f = \lambda w(\psi) f
\]
with $p = w = 1$ and $q = V_\mathrm{eff}$.  The periodic boundary conditions
are symmetric (Lagrange identity vanishes), so $\hat{H}_\psi$ is a symmetric
operator for any real, measurable $V_\mathrm{eff}$.

\textbf{Proof status}: \proved{Symmetry of $\hat{H}_\psi$ on $D_0$ for any
real $V_\mathrm{eff} \in L^2(S^1_{L_\psi})$.  Self-adjointness requires
additional work (see selfadjointness\_attempt.tex).}

%------------------------------------------------------------
\section{Pseudo-Differential and Theta-Kernel Variants}
%------------------------------------------------------------

\subsection{Fractional Laplacian variant}
\label{sec:pseudodiff}

An alternative candidate (motivated by the Berry--Keating $H = xp$ analogy)
is the pseudo-differential operator
\[
  \hat{H}^{(s)}_\psi = (-d^2/d\psi^2)^s
\]
for some $s \in (0,1)$.  Its spectrum is $\lambda_n^{(s)} = |2\pi n/L_\psi|^{2s}$,
which gives a different spectral density.

\textbf{Proof status}: \speculative{No derivation of $s \neq 1$ from UBT;
this is a phenomenological alternative.}

\subsection{Theta-kernel operator}

A trace-class candidate is the theta-kernel operator
\[
  K_t f(\psi)
  = \int_0^{L_\psi} \Theta_3\!\left(\frac{\psi-\psi'}{L_\psi}, it\right)
    f(\psi')\, d\psi',
\]
where $\Theta_3(z,\tau) = \sum_{n\in\mathbb{Z}} e^{\pi i n^2 \tau + 2\pi i nz}$
is the Jacobi theta function.  The operator $K_t$ is the heat semigroup
$e^{-t\hat{H}_\psi^{(0)}}$ for the free Hamiltonian.

\textbf{Proof status}: \proved{For the free case; the perturbed semigroup
$e^{-t\hat{H}_\psi}$ requires the Duhamel expansion and gap closure on
$V_\mathrm{eff}$ (Gap~\ref{gap:Veff}).}

%------------------------------------------------------------
\section{Open Gaps}
%------------------------------------------------------------

\begin{gap}[$L_\psi$ from UBT moduli]
  \label{gap:Lpsi}
  Derive the circumference $L_\psi$ of the compact $\psi$-direction from the
  UBT field equations and moduli space.  Currently $L_\psi = 2\pi$ is a
  conventional normalisation, not a derived quantity.

  \textbf{Assumptions required}: UBT Lagrangian fully specified including
  boundary conditions; moduli stabilisation mechanism.

  \textbf{Falsification}: If the UBT moduli space has no compact $\psi$-sector,
  the entire $\hat{H}_\psi$ programme collapses.
\end{gap}

\begin{gap}[Explicit $V_\mathrm{eff}$ derivation]
  \label{gap:Veff}
  Derive $V_\mathrm{eff}(\psi)$ explicitly from the UBT Lagrangian by
  integrating over spatial biquaternion directions.  Show that
  $V_\mathrm{eff} \in L^2(S^1_{L_\psi})$ or $V_\mathrm{eff} \in L^\infty$.

  \textbf{Assumptions required}: UBT field equations; background field
  approximation for $\langle\Theta^\dagger\Theta\rangle_q$.

  \textbf{Falsification}: If $V_\mathrm{eff}$ diverges or is not in
  $L^p$ for any $p \geq 2$, the Kato--Rellich approach fails.
\end{gap}

\begin{gap}[Connection to prime-stability $V(n)$]
  \label{gap:discrete_continuous}
  Establish the precise relation between the discrete effective potential
  $V(n) = n^2 - Bn\ln n$ (on $n \in \mathbb{N}$) and the continuous
  $V_\mathrm{eff}(\psi)$ (on $S^1_{L_\psi}$).  Is there a ``semi-classical''
  limit $\psi \leftrightarrow n$?

  \textbf{Assumptions required}: KK tower interpretation; natural units.

  \textbf{Falsification}: If the discrete and continuous spectra are structurally
  incompatible (e.g.\ different spectral dimensions), the analogy breaks.
\end{gap}

%------------------------------------------------------------
\section{Assumptions and Falsification Summary}
%------------------------------------------------------------

\begin{center}
\begin{tabular}{llll}
\hline
Claim & Status & Assumption & Falsification path \\
\hline
$\mathscr{H}_\psi = L^2(S^1)$ & \textbf{Proved} & UBT has compact $\psi$ & $\psi$ non-compact \\
Free spectrum $\lambda_n = (2\pi n)^2$ & \textbf{Proved} & $V_\mathrm{eff}=0$ & --- \\
$L_\psi = 2\pi$ & \textbf{Heuristic} & Conventional normalisation & Moduli derivation \\
$V_\mathrm{eff}$ from UBT & \textbf{Open gap} & Background field approx. & Divergent potential \\
Spectrum $\sim$ Riemann zeros & \textbf{Speculative} & G1--G6 all closed & Numerical mismatch \\
\hline
\end{tabular}
\end{center}

%------------------------------------------------------------
\section{Next Steps}
%------------------------------------------------------------

\begin{enumerate}
\item Close Gap~\ref{gap:Lpsi}: derive $L_\psi$ from UBT moduli.
\item Close Gap~\ref{gap:Veff}: compute $V_\mathrm{eff}(\psi)$ from the
      Lagrangian.
\item Verify $V_\mathrm{eff} \in L^\infty$ or $L^2$ (Kato--Rellich check).
\item Proceed to \texttt{selfadjointness\_attempt.tex}.
\item Compute the perturbed spectrum numerically for comparison against
      known Riemann zeros (see \texttt{spectral\_statistics.md}).
\end{enumerate}

%------------------------------------------------------------
\begin{thebibliography}{9}
\bibitem{BK99} M.V.\ Berry and J.P.\ Keating,
  ``$H = xp$ and the Riemann zeros,''
  \emph{Supersymmetry and Trace Formulae}, 1999.
\bibitem{C99} A.\ Connes, ``Trace formula in noncommutative geometry and
  the zeros of the Riemann zeta function,''
  \emph{Selecta Mathematica} \textbf{5} (1999), 29--106.
\bibitem{BC95} J.-B.\ Bost and A.\ Connes, ``Hecke algebras, type III
  factors and phase transitions with spontaneous symmetry breaking,''
  \emph{Selecta Mathematica} \textbf{1} (1995), 411--457.
\bibitem{ReedSimon2} M.\ Reed and B.\ Simon, \emph{Methods of Modern
  Mathematical Physics, Vol.\ II: Fourier Analysis, Self-Adjointness},
  Academic Press, 1975.
\end{thebibliography}

\end{document}
